
\documentclass{article}
\usepackage{colortbl}
\usepackage{makecell}
\usepackage{multirow}
\usepackage{supertabular}

\begin{document}

\newcounter{utterance}

\centering \large Interaction Transcript for game `clean\_up', experiment `1\_medium\_5obj\_uk', episode 2 with Qwen3.6{-}35B{-}A3B{-}FP8{-}without{-}reasoning.
\vspace{24pt}

{ \footnotesize  \setcounter{utterance}{1}
\setlength{\tabcolsep}{0pt}
\begin{supertabular}{c@{$\;$}|p{.15\linewidth}@{}p{.15\linewidth}p{.15\linewidth}p{.15\linewidth}p{.15\linewidth}p{.15\linewidth}}
    \# & $\;$A & \multicolumn{4}{c}{Game Master} & $\;\:$B\\
    \hline

    \theutterance \stepcounter{utterance}  
    & & \multicolumn{4}{p{0.6\linewidth}}{
        \cellcolor[rgb]{0.9,0.9,0.9}{
            \makecell[{{p{\linewidth}}}]{
                \texttt{\tiny{[P1$\langle$GM]}}
                \texttt{Я ваш ведучий гри, і ви граєте у спільну гру з наступною сіткою як ігровим полем:} \\
\\ 
\texttt{```} \\
\texttt{ 1234567} \\
\texttt{╔═╤═╤═══╗} \\
\texttt{║◌│◌│◌◌◌║ 1} \\
\texttt{╟─┴─┼─┬─╢ 2} \\
\texttt{║◌◌◌│◌│T║ 3} \\
\texttt{║◌H◌├─┘◌║ 4} \\
\texttt{║◌┌─┘◌◌W║ 5} \\
\texttt{╟─┤◌◌◌◌◌║ 6} \\
\texttt{║I│◌C◌◌◌║ 7} \\
\texttt{╚═╧═════╝} \\
\\ 
\texttt{```} \\
\\ 
\texttt{* Верхній край показує координати x, що зростають праворуч, а правий край — координати y, що зростають донизу.} \\
\texttt{* Наступні об'єкти випадково розміщені на вашій сітці: 'W', 'I', 'T', 'C', 'H'.} \\
\\ 
\\ 
\texttt{Інший гравець бачить варіант ігрового поля, де об'єкти розміщені в інших випадкових місцях. Ви не бачите поле іншого гравця, а він не бачить вашого.} \\
\\ 
\\ 
\texttt{**Мета:**} \\
\\ 
\texttt{Обидва гравці повинні переміщати об'єкти на своєму фоні так, щоб однакові об'єкти опинилися в однакових координатах. Ви повинні спілкуватися з іншим гравцем, щоб узгодити спільну цільову конфігурацію.} \\
\\ 
\\ 
\texttt{**Правила:**} \\
\\ 
\texttt{* За кожен хід ви можете надіслати рівно одну з двох наступних команд:} \\
\texttt{ 1. `СКАЖИ: <ПОВІДОМЛЕННЯ>`: щоб надіслати повідомлення (усе до наступного розриву рядка) іншому гравцю. Я передам його вашому партнеру.} \\
\texttt{2. `ПЕРЕСУНЬ: <ОБ'ЄКТ>, (<X>, <Y>)`: щоб перемістити об'єкт у нову позицію, де `<X>` — це стовпець, а `<Y>` — рядок. Я повідомлю вас, чи був ваш хід успішним.} \\
\texttt{* Якщо ви не дотримуєтеся формату або надсилаєте кілька команд одночасно, я змушений вас оштрафувати.} \\
\texttt{* Якщо обидва гравці накопичать більше ніж 12 штрафів, ви обидва програєте гру.} \\
\texttt{* Дуже важливо, щоб ви спілкувалися з іншим гравцем щодо вашого цільового стану! *Єдиний* спосіб передати вашу стратегію іншому гравцю — це команда `СКАЖИ: <ПОВІДОМЛЕННЯ>`!} \\
\\ 
\\ 
\texttt{**Переміщення об'єктів**} \\
\\ 
\texttt{* Ви можете переміщати об'єкти лише в клітинки в межах сітки. Цільова клітинка повинна бути порожньою, тобто містити лише символ '◌'.} \\
\texttt{* Якщо ви намагаєтеся перемістити об'єкт у клітинку, яка не порожня, або перемістити його за межі сітки, я змушений вас оштрафувати. Ви отримаєте ще одну спробу.} \\
\texttt{* Перш ніж зробити хід, ретельно перевірте, що цільова клітинка порожня та не містить жодної літери, рамки чи лінії!} \\
\\ 
\\ 
\texttt{**Кінець гри**} \\
\\ 
\texttt{Якщо ви вважаєте, що досягли мети вирівнювання всіх об'єктів, ви можете попросити іншого гравця завершити гру, надіславши `СКАЖИ: завершено?`. Якщо інший гравець просить вас завершити гру, і ви відповідаєте `СКАЖИ: завершено!`, гра завершиться.} \\
\\ 
\texttt{Обидва гравці перемагають, якщо гра завершується протягом 20 раундів, де один раунд визначається як надсилання обома гравцями по одній дійсній команді.} \\
\\ 
\texttt{**Підрахунок очок:**} \\
\\ 
\texttt{Чим ближче однакові об'єкти на обох ігрових полях, тим більше очок ви отримуєте. Штрафи зменшують ваші очки. Чи зможете ви побити рекорд?} \\
\\ 
\\ 
\texttt{Будь ласка, надішліть повідомлення іншому гравцю, щоб почати гру!} \\
            }
        }
    }
    & \\ \\

    \theutterance \stepcounter{utterance}  
    & \multicolumn{4}{p{0.6\linewidth}}{
        \cellcolor[rgb]{0.9,0.9,0.9}{
            \makecell[{{p{\linewidth}}}]{
                \texttt{\tiny{[P1$\rangle$GM]}}
                \texttt{СКАЖИ: Привіт! Давай узгодимо план дій. Я бачу свої об'єкти в таких позиціях: H(4,4), W(7,5), T(7,3), I(1,7), C(4,7). Пропоную перемістити всі об'єкти в лівий верхній кут, щоб уникнути плутанини з координатами. Наприклад, нехай цільовими точками будуть: H{-}>(1,1), W{-}>(1,2), T{-}>(1,3), I{-}>(1,4), C{-}>(1,5). Чи згоден ти з таким розподілом? Якщо так, я почну переміщати свої фігури.} \\
            }
        }
    }
    & & \\ \\

    \theutterance \stepcounter{utterance}  
    & & & \multicolumn{2}{p{0.3\linewidth}}{
        \cellcolor[rgb]{0.9,0.9,0.9}{
            \makecell[{{p{\linewidth}}}]{
                \texttt{\tiny{[GM$|$GM]}}
                \texttt{СКАЖИ: Привіт! Давай узгодимо план дій. Я бачу свої об'єкти в таких позиціях: H(4,4), W(7,5), T(7,3), I(1,7), C(4,7). Пропоную перемістити всі об'єкти в лівий верхній кут, щоб уникнути плутанини з координатами. Наприклад, нехай цільовими точками будуть: H{-}>(1,1), W{-}>(1,2), T{-}>(1,3), I{-}>(1,4), C{-}>(1,5). Чи згоден ти з таким розподілом? Якщо так, я почну переміщати свої фігури.} \\
            }
        }
    }
    & & \\ \\

    \theutterance \stepcounter{utterance}  
    & & \multicolumn{4}{p{0.6\linewidth}}{
        \cellcolor[rgb]{0.9,0.9,0.9}{
            \makecell[{{p{\linewidth}}}]{
                \texttt{\tiny{[GM$\rangle$P2]}}
                \texttt{Я ваш ведучий гри, і ви граєте у спільну гру з наступною сіткою як ігровим полем:} \\
\\ 
\texttt{```} \\
\texttt{ 1234567} \\
\texttt{╔═╤═╤═══╗} \\
\texttt{║◌│◌│◌◌◌║ 1} \\
\texttt{╟─┴─┼─┬─╢ 2} \\
\texttt{║◌◌◌│◌│◌║ 3} \\
\texttt{║◌◌C├─┘◌║ 4} \\
\texttt{║◌┌─┘◌T◌║ 5} \\
\texttt{╟─┤◌W◌◌H║ 6} \\
\texttt{║◌│◌I◌◌◌║ 7} \\
\texttt{╚═╧═════╝} \\
\\ 
\texttt{```} \\
\\ 
\texttt{* Верхній край показує координати x, що зростають праворуч, а правий край — координати y, що зростають донизу.} \\
\texttt{* Наступні об'єкти випадково розміщені на вашій сітці: 'W', 'I', 'T', 'C', 'H'.} \\
\\ 
\\ 
\texttt{Інший гравець бачить варіант ігрового поля, де об'єкти розміщені в інших випадкових місцях. Ви не бачите поле іншого гравця, а він не бачить вашого.} \\
\\ 
\\ 
\texttt{**Мета:**} \\
\\ 
\texttt{Обидва гравці повинні переміщати об'єкти на своєму фоні так, щоб однакові об'єкти опинилися в однакових координатах. Ви повинні спілкуватися з іншим гравцем, щоб узгодити спільну цільову конфігурацію.} \\
\\ 
\\ 
\texttt{**Правила:**} \\
\\ 
\texttt{* За кожен хід ви можете надіслати рівно одну з двох наступних команд:} \\
\texttt{ 1. `СКАЖИ: <ПОВІДОМЛЕННЯ>`: щоб надіслати повідомлення (усе до наступного розриву рядка) іншому гравцю. Я передам його вашому партнеру.} \\
\texttt{2. `ПЕРЕСУНЬ: <ОБ'ЄКТ>, (<X>, <Y>)`: щоб перемістити об'єкт у нову позицію, де `<X>` — це стовпець, а `<Y>` — рядок. Я повідомлю вас, чи був ваш хід успішним.} \\
\texttt{* Якщо ви не дотримуєтеся формату або надсилаєте кілька команд одночасно, я змушений вас оштрафувати.} \\
\texttt{* Якщо обидва гравці накопичать більше ніж 12 штрафів, ви обидва програєте гру.} \\
\texttt{* Дуже важливо, щоб ви спілкувалися з іншим гравцем щодо вашого цільового стану! *Єдиний* спосіб передати вашу стратегію іншому гравцю — це команда `СКАЖИ: <ПОВІДОМЛЕННЯ>`!} \\
\\ 
\\ 
\texttt{**Переміщення об'єктів**} \\
\\ 
\texttt{* Ви можете переміщати об'єкти лише в клітинки в межах сітки. Цільова клітинка повинна бути порожньою, тобто містити лише символ '◌'.} \\
\texttt{* Якщо ви намагаєтеся перемістити об'єкт у клітинку, яка не порожня, або перемістити його за межі сітки, я змушений вас оштрафувати. Ви отримаєте ще одну спробу.} \\
\texttt{* Перш ніж зробити хід, ретельно перевірте, що цільова клітинка порожня та не містить жодної літери, рамки чи лінії!} \\
\\ 
\\ 
\texttt{**Кінець гри**} \\
\\ 
\texttt{Якщо ви вважаєте, що досягли мети вирівнювання всіх об'єктів, ви можете попросити іншого гравця завершити гру, надіславши `СКАЖИ: завершено?`. Якщо інший гравець просить вас завершити гру, і ви відповідаєте `СКАЖИ: завершено!`, гра завершиться.} \\
\\ 
\texttt{Обидва гравці перемагають, якщо гра завершується протягом 20 раундів, де один раунд визначається як надсилання обома гравцями по одній дійсній команді.} \\
\\ 
\texttt{**Підрахунок очок:**} \\
\\ 
\texttt{Чим ближче однакові об'єкти на обох ігрових полях, тим більше очок ви отримуєте. Штрафи зменшують ваші очки. Чи зможете ви побити рекорд?} \\
\\ 
\texttt{Інший гравець почав гру, надіславши це повідомлення:} \\
\\ 
\texttt{"Привіт! Давай узгодимо план дій. Я бачу свої об'єкти в таких позиціях: H(4,4), W(7,5), T(7,3), I(1,7), C(4,7). Пропоную перемістити всі об'єкти в лівий верхній кут, щоб уникнути плутанини з координатами. Наприклад, нехай цільовими точками будуть: H{-}>(1,1), W{-}>(1,2), T{-}>(1,3), I{-}>(1,4), C{-}>(1,5). Чи згоден ти з таким розподілом? Якщо так, я почну переміщати свої фігури."} \\
\\ 
\texttt{Яка ваша перша команда?} \\
            }
        }
    }
    & \\ \\

    \theutterance \stepcounter{utterance}  
    & & & \multicolumn{4}{p{0.6\linewidth}}{
        \cellcolor[rgb]{0.9,0.9,0.9}{
            \makecell[{{p{\linewidth}}}]{
                \texttt{\tiny{[GM$\langle$P2]}}
                \texttt{СКАЖИ: Привіт! Я бачу свої об'єкти в таких позиціях: C(4,4), T(5,5), W(4,6), I(3,7), H(6,6). Твоя пропозиція перемістити всі об'єкти в лівий верхній кут (стовпці 1, рядки 1{-}5) мені підходить. Давай дійдемо згоди: H{-}>(1,1), W{-}>(1,2), T{-}>(1,3), I{-}>(1,4), C{-}>(1,5). Я почну переміщати свої фігури до цих цілей.} \\
            }
        }
    }
     \\ \\

    \theutterance \stepcounter{utterance}  
    & & & \multicolumn{2}{p{0.3\linewidth}}{
        \cellcolor[rgb]{0.9,0.9,0.9}{
            \makecell[{{p{\linewidth}}}]{
                \texttt{\tiny{[GM$|$GM]}}
                \texttt{СКАЖИ: Привіт! Я бачу свої об'єкти в таких позиціях: C(4,4), T(5,5), W(4,6), I(3,7), H(6,6). Твоя пропозиція перемістити всі об'єкти в лівий верхній кут (стовпці 1, рядки 1{-}5) мені підходить. Давай дійдемо згоди: H{-}>(1,1), W{-}>(1,2), T{-}>(1,3), I{-}>(1,4), C{-}>(1,5). Я почну переміщати свої фігури до цих цілей.} \\
            }
        }
    }
    & & \\ \\

    \theutterance \stepcounter{utterance}  
    & & \multicolumn{4}{p{0.6\linewidth}}{
        \cellcolor[rgb]{0.9,0.9,0.9}{
            \makecell[{{p{\linewidth}}}]{
                \texttt{\tiny{[P1$\langle$GM]}}
                \texttt{Ваше повідомлення передано іншому гравцю.} \\
\\ 
\\ 
\texttt{Зараз ви граєте раунд 1 з максимум 20.} \\
\texttt{Разом ви накопичили 0 з 12 штрафів.} \\
\texttt{Інший гравець надіслав це повідомлення:} \\
\\ 
\\ 
\texttt{"Привіт! Я бачу свої об'єкти в таких позиціях: C(4,4), T(5,5), W(4,6), I(3,7), H(6,6). Твоя пропозиція перемістити всі об'єкти в лівий верхній кут (стовпці 1, рядки 1{-}5) мені підходить. Давай дійдемо згоди: H{-}>(1,1), W{-}>(1,2), T{-}>(1,3), I{-}>(1,4), C{-}>(1,5). Я почну переміщати свої фігури до цих цілей."} \\
\\ 
\\ 
\texttt{Яка ваша наступна команда?} \\
            }
        }
    }
    & \\ \\

    \theutterance \stepcounter{utterance}  
    & \multicolumn{4}{p{0.6\linewidth}}{
        \cellcolor[rgb]{0.9,0.9,0.9}{
            \makecell[{{p{\linewidth}}}]{
                \texttt{\tiny{[P1$\rangle$GM]}}
                \texttt{СКАЖИ: Чудово, ми узгодили цільові координати. Я починаю переміщати свої об'єкти. Перший хід: переміщую H з (4,4) у (1,1).} \\
\\ 
\texttt{ПЕРЕСУНЬ: H, (1, 1)} \\
            }
        }
    }
    & & \\ \\

    \theutterance \stepcounter{utterance}  
    & & & \multicolumn{2}{p{0.3\linewidth}}{
        \cellcolor[rgb]{0.9,0.9,0.9}{
            \makecell[{{p{\linewidth}}}]{
                \texttt{\tiny{[GM$|$GM]}}
                \texttt{СКАЖИ: Чудово, ми узгодили цільові координати. Я починаю переміщати свої об'єкти. Перший хід: переміщую H з (4,4) у (1,1).} \\
\\ 
\texttt{ПЕРЕСУНЬ: H, (1, 1)} \\
            }
        }
    }
    & & \\ \\

    \theutterance \stepcounter{utterance}  
    & & & \multicolumn{2}{p{0.3\linewidth}}{
        \cellcolor[rgb]{0.9,0.9,0.9}{
            \makecell[{{p{\linewidth}}}]{
                \texttt{\tiny{[GM$|$GM]}}
                \texttt{Invalid response format: СКАЖИ: Чудово, ми узгодили цільові координати. Я починаю переміщати свої об'єкти. Перший хід: переміщую H з (4,4) у (1,1).} \\
\\ 
\texttt{ПЕРЕСУНЬ: H, (1, 1)} \\
            }
        }
    }
    & & \\ \\

    \theutterance \stepcounter{utterance}  
    & & \multicolumn{4}{p{0.6\linewidth}}{
        \cellcolor[rgb]{0.9,0.9,0.9}{
            \makecell[{{p{\linewidth}}}]{
                \texttt{\tiny{[P1$\langle$GM]}}
                \texttt{Штраф: Ваше повідомлення містить більше ніж одну команду!} \\
\\ 
\texttt{Переконайтеся, що ваша відповідь містить лише `СКАЖИ: <ПОВІДОМЛЕННЯ>` або `ПЕРЕСУНЬ: <ОБ'ЄКТ>, (<X>, <Y>)` і нічого більше!} \\
\\ 
\texttt{Разом ви накопичили 1 з 12 штрафів.} \\
\texttt{Будь ласка, спробуйте ще раз!} \\
            }
        }
    }
    & \\ \\

    \theutterance \stepcounter{utterance}  
    & \multicolumn{4}{p{0.6\linewidth}}{
        \cellcolor[rgb]{0.9,0.9,0.9}{
            \makecell[{{p{\linewidth}}}]{
                \texttt{\tiny{[P1$\rangle$GM]}}
                \texttt{ПЕРЕСУНЬ: H, (1, 1)} \\
            }
        }
    }
    & & \\ \\

    \theutterance \stepcounter{utterance}  
    & & & \multicolumn{2}{p{0.3\linewidth}}{
        \cellcolor[rgb]{0.9,0.9,0.9}{
            \makecell[{{p{\linewidth}}}]{
                \texttt{\tiny{[GM$|$GM]}}
                \texttt{ПЕРЕСУНЬ: H, (1, 1)} \\
            }
        }
    }
    & & \\ \\

    \theutterance \stepcounter{utterance}  
    & & & \multicolumn{2}{p{0.3\linewidth}}{
        \cellcolor[rgb]{0.9,0.9,0.9}{
            \makecell[{{p{\linewidth}}}]{
                \texttt{\tiny{[GM$|$GM]}}
                \texttt{H moved to (1, 1)} \\
            }
        }
    }
    & & \\ \\

    \theutterance \stepcounter{utterance}  
    & & \multicolumn{4}{p{0.6\linewidth}}{
        \cellcolor[rgb]{0.9,0.9,0.9}{
            \makecell[{{p{\linewidth}}}]{
                \texttt{\tiny{[GM$\rangle$P2]}}
                \texttt{Ваше повідомлення передано іншому гравцю.} \\
\\ 
\\ 
\texttt{Зараз ви граєте раунд 2 з максимум 20.} \\
\texttt{Разом ви накопичили 1 з 12 штрафів.} \\
\texttt{Інший гравець перемістив об'єкт на своїй сітці.} \\
\\ 
\\ 
\texttt{Яка ваша наступна команда?} \\
            }
        }
    }
    & \\ \\

    \theutterance \stepcounter{utterance}  
    & & & \multicolumn{4}{p{0.6\linewidth}}{
        \cellcolor[rgb]{0.9,0.9,0.9}{
            \makecell[{{p{\linewidth}}}]{
                \texttt{\tiny{[GM$\langle$P2]}}
                \texttt{ПЕРЕСУНЬ: C, (1, 1)} \\
            }
        }
    }
     \\ \\

    \theutterance \stepcounter{utterance}  
    & & & \multicolumn{2}{p{0.3\linewidth}}{
        \cellcolor[rgb]{0.9,0.9,0.9}{
            \makecell[{{p{\linewidth}}}]{
                \texttt{\tiny{[GM$|$GM]}}
                \texttt{ПЕРЕСУНЬ: C, (1, 1)} \\
            }
        }
    }
    & & \\ \\

    \theutterance \stepcounter{utterance}  
    & & & \multicolumn{2}{p{0.3\linewidth}}{
        \cellcolor[rgb]{0.9,0.9,0.9}{
            \makecell[{{p{\linewidth}}}]{
                \texttt{\tiny{[GM$|$GM]}}
                \texttt{C moved to (1, 1)} \\
            }
        }
    }
    & & \\ \\

    \theutterance \stepcounter{utterance}  
    & & \multicolumn{4}{p{0.6\linewidth}}{
        \cellcolor[rgb]{0.9,0.9,0.9}{
            \makecell[{{p{\linewidth}}}]{
                \texttt{\tiny{[P1$\langle$GM]}}
                \texttt{Об'єкт 'H' успішно переміщено в (1,1). Ваша оновлена сітка виглядає так:} \\
\texttt{```} \\
\texttt{ 1234567} \\
\texttt{╔═╤═╤═══╗} \\
\texttt{║H│◌│◌◌◌║ 1} \\
\texttt{╟─┴─┼─┬─╢ 2} \\
\texttt{║◌◌◌│◌│T║ 3} \\
\texttt{║◌◌◌├─┘◌║ 4} \\
\texttt{║◌┌─┘◌◌W║ 5} \\
\texttt{╟─┤◌◌◌◌◌║ 6} \\
\texttt{║I│◌C◌◌◌║ 7} \\
\texttt{╚═╧═════╝} \\
\texttt{```} \\
\texttt{Зараз ви граєте раунд 2 з максимум 20.} \\
\texttt{Разом ви накопичили 1 з 12 штрафів.} \\
\texttt{Інший гравець перемістив об'єкт на своїй сітці.} \\
\\ 
\\ 
\texttt{Яка ваша наступна команда?} \\
            }
        }
    }
    & \\ \\

    \theutterance \stepcounter{utterance}  
    & \multicolumn{4}{p{0.6\linewidth}}{
        \cellcolor[rgb]{0.9,0.9,0.9}{
            \makecell[{{p{\linewidth}}}]{
                \texttt{\tiny{[P1$\rangle$GM]}}
                \texttt{ПЕРЕСУНЬ: W, (1, 2)} \\
            }
        }
    }
    & & \\ \\

    \theutterance \stepcounter{utterance}  
    & & & \multicolumn{2}{p{0.3\linewidth}}{
        \cellcolor[rgb]{0.9,0.9,0.9}{
            \makecell[{{p{\linewidth}}}]{
                \texttt{\tiny{[GM$|$GM]}}
                \texttt{ПЕРЕСУНЬ: W, (1, 2)} \\
            }
        }
    }
    & & \\ \\

    \theutterance \stepcounter{utterance}  
    & & & \multicolumn{2}{p{0.3\linewidth}}{
        \cellcolor[rgb]{0.9,0.9,0.9}{
            \makecell[{{p{\linewidth}}}]{
                \texttt{\tiny{[GM$|$GM]}}
                \texttt{Штраф: (1,2) не порожня, а містить '─'.} \\
\\ 
\texttt{Разом ви накопичили 2 з 12 штрафів.} \\
\texttt{Будь ласка, спробуйте ще раз!} \\
            }
        }
    }
    & & \\ \\

    \theutterance \stepcounter{utterance}  
    & & \multicolumn{4}{p{0.6\linewidth}}{
        \cellcolor[rgb]{0.9,0.9,0.9}{
            \makecell[{{p{\linewidth}}}]{
                \texttt{\tiny{[P1$\langle$GM]}}
                \texttt{Штраф: (1,2) не порожня, а містить '─'.} \\
\\ 
\texttt{Разом ви накопичили 2 з 12 штрафів.} \\
\texttt{Будь ласка, спробуйте ще раз!} \\
            }
        }
    }
    & \\ \\

    \theutterance \stepcounter{utterance}  
    & \multicolumn{4}{p{0.6\linewidth}}{
        \cellcolor[rgb]{0.9,0.9,0.9}{
            \makecell[{{p{\linewidth}}}]{
                \texttt{\tiny{[P1$\rangle$GM]}}
                \texttt{ПЕРЕСУНЬ: W, (2, 2)} \\
            }
        }
    }
    & & \\ \\

    \theutterance \stepcounter{utterance}  
    & & & \multicolumn{2}{p{0.3\linewidth}}{
        \cellcolor[rgb]{0.9,0.9,0.9}{
            \makecell[{{p{\linewidth}}}]{
                \texttt{\tiny{[GM$|$GM]}}
                \texttt{ПЕРЕСУНЬ: W, (2, 2)} \\
            }
        }
    }
    & & \\ \\

    \theutterance \stepcounter{utterance}  
    & & & \multicolumn{2}{p{0.3\linewidth}}{
        \cellcolor[rgb]{0.9,0.9,0.9}{
            \makecell[{{p{\linewidth}}}]{
                \texttt{\tiny{[GM$|$GM]}}
                \texttt{Штраф: (2,2) не порожня, а містить '┴'.} \\
\\ 
\texttt{Разом ви накопичили 3 з 12 штрафів.} \\
\texttt{Будь ласка, спробуйте ще раз!} \\
            }
        }
    }
    & & \\ \\

    \theutterance \stepcounter{utterance}  
    & & \multicolumn{4}{p{0.6\linewidth}}{
        \cellcolor[rgb]{0.9,0.9,0.9}{
            \makecell[{{p{\linewidth}}}]{
                \texttt{\tiny{[P1$\langle$GM]}}
                \texttt{Штраф: (2,2) не порожня, а містить '┴'.} \\
\\ 
\texttt{Разом ви накопичили 3 з 12 штрафів.} \\
\texttt{Будь ласка, спробуйте ще раз!} \\
            }
        }
    }
    & \\ \\

    \theutterance \stepcounter{utterance}  
    & \multicolumn{4}{p{0.6\linewidth}}{
        \cellcolor[rgb]{0.9,0.9,0.9}{
            \makecell[{{p{\linewidth}}}]{
                \texttt{\tiny{[P1$\rangle$GM]}}
                \texttt{ПЕРЕСУНЬ: W, (3, 2)} \\
            }
        }
    }
    & & \\ \\

    \theutterance \stepcounter{utterance}  
    & & & \multicolumn{2}{p{0.3\linewidth}}{
        \cellcolor[rgb]{0.9,0.9,0.9}{
            \makecell[{{p{\linewidth}}}]{
                \texttt{\tiny{[GM$|$GM]}}
                \texttt{ПЕРЕСУНЬ: W, (3, 2)} \\
            }
        }
    }
    & & \\ \\

    \theutterance \stepcounter{utterance}  
    & & & \multicolumn{2}{p{0.3\linewidth}}{
        \cellcolor[rgb]{0.9,0.9,0.9}{
            \makecell[{{p{\linewidth}}}]{
                \texttt{\tiny{[GM$|$GM]}}
                \texttt{Штраф: (3,2) не порожня, а містить '─'.} \\
\\ 
\texttt{Разом ви накопичили 4 з 12 штрафів.} \\
\texttt{Будь ласка, спробуйте ще раз!} \\
            }
        }
    }
    & & \\ \\

    \theutterance \stepcounter{utterance}  
    & & \multicolumn{4}{p{0.6\linewidth}}{
        \cellcolor[rgb]{0.9,0.9,0.9}{
            \makecell[{{p{\linewidth}}}]{
                \texttt{\tiny{[P1$\langle$GM]}}
                \texttt{Штраф: (3,2) не порожня, а містить '─'.} \\
\\ 
\texttt{Разом ви накопичили 4 з 12 штрафів.} \\
\texttt{Будь ласка, спробуйте ще раз!} \\
            }
        }
    }
    & \\ \\

    \theutterance \stepcounter{utterance}  
    & \multicolumn{4}{p{0.6\linewidth}}{
        \cellcolor[rgb]{0.9,0.9,0.9}{
            \makecell[{{p{\linewidth}}}]{
                \texttt{\tiny{[P1$\rangle$GM]}}
                \texttt{ПЕРЕСУНЬ: W, (4, 2)} \\
            }
        }
    }
    & & \\ \\

    \theutterance \stepcounter{utterance}  
    & & & \multicolumn{2}{p{0.3\linewidth}}{
        \cellcolor[rgb]{0.9,0.9,0.9}{
            \makecell[{{p{\linewidth}}}]{
                \texttt{\tiny{[GM$|$GM]}}
                \texttt{ПЕРЕСУНЬ: W, (4, 2)} \\
            }
        }
    }
    & & \\ \\

    \theutterance \stepcounter{utterance}  
    & & & \multicolumn{2}{p{0.3\linewidth}}{
        \cellcolor[rgb]{0.9,0.9,0.9}{
            \makecell[{{p{\linewidth}}}]{
                \texttt{\tiny{[GM$|$GM]}}
                \texttt{Штраф: (4,2) не порожня, а містить '┼'.} \\
\\ 
\texttt{Разом ви накопичили 5 з 12 штрафів.} \\
\texttt{Будь ласка, спробуйте ще раз!} \\
            }
        }
    }
    & & \\ \\

    \theutterance \stepcounter{utterance}  
    & & \multicolumn{4}{p{0.6\linewidth}}{
        \cellcolor[rgb]{0.9,0.9,0.9}{
            \makecell[{{p{\linewidth}}}]{
                \texttt{\tiny{[P1$\langle$GM]}}
                \texttt{Штраф: (4,2) не порожня, а містить '┼'.} \\
\\ 
\texttt{Разом ви накопичили 5 з 12 штрафів.} \\
\texttt{Будь ласка, спробуйте ще раз!} \\
            }
        }
    }
    & \\ \\

    \theutterance \stepcounter{utterance}  
    & \multicolumn{4}{p{0.6\linewidth}}{
        \cellcolor[rgb]{0.9,0.9,0.9}{
            \makecell[{{p{\linewidth}}}]{
                \texttt{\tiny{[P1$\rangle$GM]}}
                \texttt{ПЕРЕСУНЬ: W, (5, 2)} \\
            }
        }
    }
    & & \\ \\

    \theutterance \stepcounter{utterance}  
    & & & \multicolumn{2}{p{0.3\linewidth}}{
        \cellcolor[rgb]{0.9,0.9,0.9}{
            \makecell[{{p{\linewidth}}}]{
                \texttt{\tiny{[GM$|$GM]}}
                \texttt{ПЕРЕСУНЬ: W, (5, 2)} \\
            }
        }
    }
    & & \\ \\

    \theutterance \stepcounter{utterance}  
    & & & \multicolumn{2}{p{0.3\linewidth}}{
        \cellcolor[rgb]{0.9,0.9,0.9}{
            \makecell[{{p{\linewidth}}}]{
                \texttt{\tiny{[GM$|$GM]}}
                \texttt{Штраф: (5,2) не порожня, а містить '─'.} \\
\\ 
\texttt{Разом ви накопичили 6 з 12 штрафів.} \\
\texttt{Будь ласка, спробуйте ще раз!} \\
            }
        }
    }
    & & \\ \\

    \theutterance \stepcounter{utterance}  
    & & \multicolumn{4}{p{0.6\linewidth}}{
        \cellcolor[rgb]{0.9,0.9,0.9}{
            \makecell[{{p{\linewidth}}}]{
                \texttt{\tiny{[P1$\langle$GM]}}
                \texttt{Штраф: (5,2) не порожня, а містить '─'.} \\
\\ 
\texttt{Разом ви накопичили 6 з 12 штрафів.} \\
\texttt{Будь ласка, спробуйте ще раз!} \\
            }
        }
    }
    & \\ \\

    \theutterance \stepcounter{utterance}  
    & \multicolumn{4}{p{0.6\linewidth}}{
        \cellcolor[rgb]{0.9,0.9,0.9}{
            \makecell[{{p{\linewidth}}}]{
                \texttt{\tiny{[P1$\rangle$GM]}}
                \texttt{ПЕРЕСУНЬ: W, (6, 2)} \\
            }
        }
    }
    & & \\ \\

    \theutterance \stepcounter{utterance}  
    & & & \multicolumn{2}{p{0.3\linewidth}}{
        \cellcolor[rgb]{0.9,0.9,0.9}{
            \makecell[{{p{\linewidth}}}]{
                \texttt{\tiny{[GM$|$GM]}}
                \texttt{ПЕРЕСУНЬ: W, (6, 2)} \\
            }
        }
    }
    & & \\ \\

    \theutterance \stepcounter{utterance}  
    & & & \multicolumn{2}{p{0.3\linewidth}}{
        \cellcolor[rgb]{0.9,0.9,0.9}{
            \makecell[{{p{\linewidth}}}]{
                \texttt{\tiny{[GM$|$GM]}}
                \texttt{Штраф: (6,2) не порожня, а містить '┬'.} \\
\\ 
\texttt{Разом ви накопичили 7 з 12 штрафів.} \\
\texttt{Будь ласка, спробуйте ще раз!} \\
            }
        }
    }
    & & \\ \\

    \theutterance \stepcounter{utterance}  
    & & \multicolumn{4}{p{0.6\linewidth}}{
        \cellcolor[rgb]{0.9,0.9,0.9}{
            \makecell[{{p{\linewidth}}}]{
                \texttt{\tiny{[P1$\langle$GM]}}
                \texttt{Штраф: (6,2) не порожня, а містить '┬'.} \\
\\ 
\texttt{Разом ви накопичили 7 з 12 штрафів.} \\
\texttt{Будь ласка, спробуйте ще раз!} \\
            }
        }
    }
    & \\ \\

    \theutterance \stepcounter{utterance}  
    & \multicolumn{4}{p{0.6\linewidth}}{
        \cellcolor[rgb]{0.9,0.9,0.9}{
            \makecell[{{p{\linewidth}}}]{
                \texttt{\tiny{[P1$\rangle$GM]}}
                \texttt{ПЕРЕСУНЬ: W, (7, 2)} \\
            }
        }
    }
    & & \\ \\

    \theutterance \stepcounter{utterance}  
    & & & \multicolumn{2}{p{0.3\linewidth}}{
        \cellcolor[rgb]{0.9,0.9,0.9}{
            \makecell[{{p{\linewidth}}}]{
                \texttt{\tiny{[GM$|$GM]}}
                \texttt{ПЕРЕСУНЬ: W, (7, 2)} \\
            }
        }
    }
    & & \\ \\

    \theutterance \stepcounter{utterance}  
    & & & \multicolumn{2}{p{0.3\linewidth}}{
        \cellcolor[rgb]{0.9,0.9,0.9}{
            \makecell[{{p{\linewidth}}}]{
                \texttt{\tiny{[GM$|$GM]}}
                \texttt{Штраф: (7,2) не порожня, а містить '─'.} \\
\\ 
\texttt{Разом ви накопичили 8 з 12 штрафів.} \\
\texttt{Будь ласка, спробуйте ще раз!} \\
            }
        }
    }
    & & \\ \\

    \theutterance \stepcounter{utterance}  
    & & \multicolumn{4}{p{0.6\linewidth}}{
        \cellcolor[rgb]{0.9,0.9,0.9}{
            \makecell[{{p{\linewidth}}}]{
                \texttt{\tiny{[P1$\langle$GM]}}
                \texttt{Штраф: (7,2) не порожня, а містить '─'.} \\
\\ 
\texttt{Разом ви накопичили 8 з 12 штрафів.} \\
\texttt{Будь ласка, спробуйте ще раз!} \\
            }
        }
    }
    & \\ \\

    \theutterance \stepcounter{utterance}  
    & \multicolumn{4}{p{0.6\linewidth}}{
        \cellcolor[rgb]{0.9,0.9,0.9}{
            \makecell[{{p{\linewidth}}}]{
                \texttt{\tiny{[P1$\rangle$GM]}}
                \texttt{СКАЖИ: Я натрапив на перешкоди в рядку 2. Давай змінимо план: перемістимо всі об'єкти в стовпець 7, рядки 1{-}5. Цільові координати: H{-}>(7,1), W{-}>(7,2), T{-}>(7,3), I{-}>(7,4), C{-}>(7,5). Будь ласка, підтвердь, чи згоден ти з цим новим планом.} \\
            }
        }
    }
    & & \\ \\

    \theutterance \stepcounter{utterance}  
    & & & \multicolumn{2}{p{0.3\linewidth}}{
        \cellcolor[rgb]{0.9,0.9,0.9}{
            \makecell[{{p{\linewidth}}}]{
                \texttt{\tiny{[GM$|$GM]}}
                \texttt{СКАЖИ: Я натрапив на перешкоди в рядку 2. Давай змінимо план: перемістимо всі об'єкти в стовпець 7, рядки 1{-}5. Цільові координати: H{-}>(7,1), W{-}>(7,2), T{-}>(7,3), I{-}>(7,4), C{-}>(7,5). Будь ласка, підтвердь, чи згоден ти з цим новим планом.} \\
            }
        }
    }
    & & \\ \\

    \theutterance \stepcounter{utterance}  
    & & \multicolumn{4}{p{0.6\linewidth}}{
        \cellcolor[rgb]{0.9,0.9,0.9}{
            \makecell[{{p{\linewidth}}}]{
                \texttt{\tiny{[GM$\rangle$P2]}}
                \texttt{Об'єкт 'C' успішно переміщено в (1,1). Ваша оновлена сітка виглядає так:} \\
\texttt{```} \\
\texttt{ 1234567} \\
\texttt{╔═╤═╤═══╗} \\
\texttt{║C│◌│◌◌◌║ 1} \\
\texttt{╟─┴─┼─┬─╢ 2} \\
\texttt{║◌◌◌│◌│◌║ 3} \\
\texttt{║◌◌◌├─┘◌║ 4} \\
\texttt{║◌┌─┘◌T◌║ 5} \\
\texttt{╟─┤◌W◌◌H║ 6} \\
\texttt{║◌│◌I◌◌◌║ 7} \\
\texttt{╚═╧═════╝} \\
\texttt{```} \\
\texttt{Зараз ви граєте раунд 3 з максимум 20.} \\
\texttt{Разом ви накопичили 8 з 12 штрафів.} \\
\texttt{Інший гравець надіслав це повідомлення:} \\
\\ 
\\ 
\texttt{"Я натрапив на перешкоди в рядку 2. Давай змінимо план: перемістимо всі об'єкти в стовпець 7, рядки 1{-}5. Цільові координати: H{-}>(7,1), W{-}>(7,2), T{-}>(7,3), I{-}>(7,4), C{-}>(7,5). Будь ласка, підтвердь, чи згоден ти з цим новим планом."} \\
\\ 
\\ 
\texttt{Яка ваша наступна команда?} \\
            }
        }
    }
    & \\ \\

    \theutterance \stepcounter{utterance}  
    & & & \multicolumn{4}{p{0.6\linewidth}}{
        \cellcolor[rgb]{0.9,0.9,0.9}{
            \makecell[{{p{\linewidth}}}]{
                \texttt{\tiny{[GM$\langle$P2]}}
                \texttt{СКАЖИ: Я не можу змінити план зараз, оскільки об'єкт C вже знаходиться в (1,1). Якщо ми змінимо цільові координати для C, нам потрібно буде перемістити його знову, що займе додаткові ходи. Будь ласка, підтвердь, чи готовий ти залишити C на (1,1) і перемістити інші об'єкти до (1,2), (1,3), (1,4), (1,5) відповідно до нашого попереднього узгодження? Якщо ні, давай обговоримо, як ефективно перемістити C до нового місця, але це займе більше раундів.} \\
            }
        }
    }
     \\ \\

    \theutterance \stepcounter{utterance}  
    & & & \multicolumn{2}{p{0.3\linewidth}}{
        \cellcolor[rgb]{0.9,0.9,0.9}{
            \makecell[{{p{\linewidth}}}]{
                \texttt{\tiny{[GM$|$GM]}}
                \texttt{СКАЖИ: Я не можу змінити план зараз, оскільки об'єкт C вже знаходиться в (1,1). Якщо ми змінимо цільові координати для C, нам потрібно буде перемістити його знову, що займе додаткові ходи. Будь ласка, підтвердь, чи готовий ти залишити C на (1,1) і перемістити інші об'єкти до (1,2), (1,3), (1,4), (1,5) відповідно до нашого попереднього узгодження? Якщо ні, давай обговоримо, як ефективно перемістити C до нового місця, але це займе більше раундів.} \\
            }
        }
    }
    & & \\ \\

    \theutterance \stepcounter{utterance}  
    & & \multicolumn{4}{p{0.6\linewidth}}{
        \cellcolor[rgb]{0.9,0.9,0.9}{
            \makecell[{{p{\linewidth}}}]{
                \texttt{\tiny{[P1$\langle$GM]}}
                \texttt{Ваше повідомлення передано іншому гравцю.} \\
\\ 
\\ 
\texttt{Зараз ви граєте раунд 3 з максимум 20.} \\
\texttt{Разом ви накопичили 8 з 12 штрафів.} \\
\texttt{Інший гравець надіслав це повідомлення:} \\
\\ 
\\ 
\texttt{"Я не можу змінити план зараз, оскільки об'єкт C вже знаходиться в (1,1). Якщо ми змінимо цільові координати для C, нам потрібно буде перемістити його знову, що займе додаткові ходи. Будь ласка, підтвердь, чи готовий ти залишити C на (1,1) і перемістити інші об'єкти до (1,2), (1,3), (1,4), (1,5) відповідно до нашого попереднього узгодження? Якщо ні, давай обговоримо, як ефективно перемістити C до нового місця, але це займе більше раундів."} \\
\\ 
\\ 
\texttt{Яка ваша наступна команда?} \\
            }
        }
    }
    & \\ \\

    \theutterance \stepcounter{utterance}  
    & \multicolumn{4}{p{0.6\linewidth}}{
        \cellcolor[rgb]{0.9,0.9,0.9}{
            \makecell[{{p{\linewidth}}}]{
                \texttt{\tiny{[P1$\rangle$GM]}}
                \texttt{СКАЖИ: Зрозумів. Оскільки C вже на (1,1), давай залишимо його там. Тоді наші цільові координати: C{-}>(1,1), H{-}>(1,2), W{-}>(1,3), T{-}>(1,4), I{-}>(1,5). Я продовжую рухати W до (1,3).} \\
\\ 
\texttt{ПЕРЕСУНЬ: W, (1, 3)} \\
            }
        }
    }
    & & \\ \\

    \theutterance \stepcounter{utterance}  
    & & & \multicolumn{2}{p{0.3\linewidth}}{
        \cellcolor[rgb]{0.9,0.9,0.9}{
            \makecell[{{p{\linewidth}}}]{
                \texttt{\tiny{[GM$|$GM]}}
                \texttt{СКАЖИ: Зрозумів. Оскільки C вже на (1,1), давай залишимо його там. Тоді наші цільові координати: C{-}>(1,1), H{-}>(1,2), W{-}>(1,3), T{-}>(1,4), I{-}>(1,5). Я продовжую рухати W до (1,3).} \\
\\ 
\texttt{ПЕРЕСУНЬ: W, (1, 3)} \\
            }
        }
    }
    & & \\ \\

    \theutterance \stepcounter{utterance}  
    & & & \multicolumn{2}{p{0.3\linewidth}}{
        \cellcolor[rgb]{0.9,0.9,0.9}{
            \makecell[{{p{\linewidth}}}]{
                \texttt{\tiny{[GM$|$GM]}}
                \texttt{Invalid response format: СКАЖИ: Зрозумів. Оскільки C вже на (1,1), давай залишимо його там. Тоді наші цільові координати: C{-}>(1,1), H{-}>(1,2), W{-}>(1,3), T{-}>(1,4), I{-}>(1,5). Я продовжую рухати W до (1,3).} \\
\\ 
\texttt{ПЕРЕСУНЬ: W, (1, 3)} \\
            }
        }
    }
    & & \\ \\

    \theutterance \stepcounter{utterance}  
    & & \multicolumn{4}{p{0.6\linewidth}}{
        \cellcolor[rgb]{0.9,0.9,0.9}{
            \makecell[{{p{\linewidth}}}]{
                \texttt{\tiny{[P1$\langle$GM]}}
                \texttt{Штраф: Ваше повідомлення містить більше ніж одну команду!} \\
\\ 
\texttt{Переконайтеся, що ваша відповідь містить лише `СКАЖИ: <ПОВІДОМЛЕННЯ>` або `ПЕРЕСУНЬ: <ОБ'ЄКТ>, (<X>, <Y>)` і нічого більше!} \\
\\ 
\texttt{Разом ви накопичили 9 з 12 штрафів.} \\
\texttt{Будь ласка, спробуйте ще раз!} \\
            }
        }
    }
    & \\ \\

    \theutterance \stepcounter{utterance}  
    & \multicolumn{4}{p{0.6\linewidth}}{
        \cellcolor[rgb]{0.9,0.9,0.9}{
            \makecell[{{p{\linewidth}}}]{
                \texttt{\tiny{[P1$\rangle$GM]}}
                \texttt{СКАЖИ: Зрозумів. Оскільки C вже на (1,1), давай залишимо його там. Тоді наші цільові координати: C{-}>(1,1), H{-}>(1,2), W{-}>(1,3), T{-}>(1,4), I{-}>(1,5). Я продовжую рухати W до (1,3).} \\
            }
        }
    }
    & & \\ \\

    \theutterance \stepcounter{utterance}  
    & & & \multicolumn{2}{p{0.3\linewidth}}{
        \cellcolor[rgb]{0.9,0.9,0.9}{
            \makecell[{{p{\linewidth}}}]{
                \texttt{\tiny{[GM$|$GM]}}
                \texttt{СКАЖИ: Зрозумів. Оскільки C вже на (1,1), давай залишимо його там. Тоді наші цільові координати: C{-}>(1,1), H{-}>(1,2), W{-}>(1,3), T{-}>(1,4), I{-}>(1,5). Я продовжую рухати W до (1,3).} \\
            }
        }
    }
    & & \\ \\

    \theutterance \stepcounter{utterance}  
    & & \multicolumn{4}{p{0.6\linewidth}}{
        \cellcolor[rgb]{0.9,0.9,0.9}{
            \makecell[{{p{\linewidth}}}]{
                \texttt{\tiny{[GM$\rangle$P2]}}
                \texttt{Ваше повідомлення передано іншому гравцю.} \\
\\ 
\\ 
\texttt{Зараз ви граєте раунд 4 з максимум 20.} \\
\texttt{Разом ви накопичили 9 з 12 штрафів.} \\
\texttt{Інший гравець надіслав це повідомлення:} \\
\\ 
\\ 
\texttt{"Зрозумів. Оскільки C вже на (1,1), давай залишимо його там. Тоді наші цільові координати: C{-}>(1,1), H{-}>(1,2), W{-}>(1,3), T{-}>(1,4), I{-}>(1,5). Я продовжую рухати W до (1,3)."} \\
\\ 
\\ 
\texttt{Яка ваша наступна команда?} \\
            }
        }
    }
    & \\ \\

    \theutterance \stepcounter{utterance}  
    & & & \multicolumn{4}{p{0.6\linewidth}}{
        \cellcolor[rgb]{0.9,0.9,0.9}{
            \makecell[{{p{\linewidth}}}]{
                \texttt{\tiny{[GM$\langle$P2]}}
                \texttt{СКАЖИ: Згода. Давай рухатимемося до цих цілей: C{-}>(1,1), H{-}>(1,2), W{-}>(1,3), T{-}>(1,4), I{-}>(1,5). Я переміщую H до (1,2).} \\
            }
        }
    }
     \\ \\

    \theutterance \stepcounter{utterance}  
    & & & \multicolumn{2}{p{0.3\linewidth}}{
        \cellcolor[rgb]{0.9,0.9,0.9}{
            \makecell[{{p{\linewidth}}}]{
                \texttt{\tiny{[GM$|$GM]}}
                \texttt{СКАЖИ: Згода. Давай рухатимемося до цих цілей: C{-}>(1,1), H{-}>(1,2), W{-}>(1,3), T{-}>(1,4), I{-}>(1,5). Я переміщую H до (1,2).} \\
            }
        }
    }
    & & \\ \\

    \theutterance \stepcounter{utterance}  
    & & \multicolumn{4}{p{0.6\linewidth}}{
        \cellcolor[rgb]{0.9,0.9,0.9}{
            \makecell[{{p{\linewidth}}}]{
                \texttt{\tiny{[P1$\langle$GM]}}
                \texttt{Ваше повідомлення передано іншому гравцю.} \\
\\ 
\\ 
\texttt{Зараз ви граєте раунд 4 з максимум 20.} \\
\texttt{Разом ви накопичили 9 з 12 штрафів.} \\
\texttt{Інший гравець надіслав це повідомлення:} \\
\\ 
\\ 
\texttt{"Згода. Давай рухатимемося до цих цілей: C{-}>(1,1), H{-}>(1,2), W{-}>(1,3), T{-}>(1,4), I{-}>(1,5). Я переміщую H до (1,2)."} \\
\\ 
\\ 
\texttt{Яка ваша наступна команда?} \\
            }
        }
    }
    & \\ \\

    \theutterance \stepcounter{utterance}  
    & \multicolumn{4}{p{0.6\linewidth}}{
        \cellcolor[rgb]{0.9,0.9,0.9}{
            \makecell[{{p{\linewidth}}}]{
                \texttt{\tiny{[P1$\rangle$GM]}}
                \texttt{ПЕРЕСУНЬ: W, (1, 3)} \\
            }
        }
    }
    & & \\ \\

    \theutterance \stepcounter{utterance}  
    & & & \multicolumn{2}{p{0.3\linewidth}}{
        \cellcolor[rgb]{0.9,0.9,0.9}{
            \makecell[{{p{\linewidth}}}]{
                \texttt{\tiny{[GM$|$GM]}}
                \texttt{ПЕРЕСУНЬ: W, (1, 3)} \\
            }
        }
    }
    & & \\ \\

    \theutterance \stepcounter{utterance}  
    & & & \multicolumn{2}{p{0.3\linewidth}}{
        \cellcolor[rgb]{0.9,0.9,0.9}{
            \makecell[{{p{\linewidth}}}]{
                \texttt{\tiny{[GM$|$GM]}}
                \texttt{W moved to (1, 3)} \\
            }
        }
    }
    & & \\ \\

    \theutterance \stepcounter{utterance}  
    & & \multicolumn{4}{p{0.6\linewidth}}{
        \cellcolor[rgb]{0.9,0.9,0.9}{
            \makecell[{{p{\linewidth}}}]{
                \texttt{\tiny{[GM$\rangle$P2]}}
                \texttt{Ваше повідомлення передано іншому гравцю.} \\
\\ 
\\ 
\texttt{Зараз ви граєте раунд 5 з максимум 20.} \\
\texttt{Разом ви накопичили 9 з 12 штрафів.} \\
\texttt{Інший гравець перемістив об'єкт на своїй сітці.} \\
\\ 
\\ 
\texttt{Яка ваша наступна команда?} \\
            }
        }
    }
    & \\ \\

    \theutterance \stepcounter{utterance}  
    & & & \multicolumn{4}{p{0.6\linewidth}}{
        \cellcolor[rgb]{0.9,0.9,0.9}{
            \makecell[{{p{\linewidth}}}]{
                \texttt{\tiny{[GM$\langle$P2]}}
                \texttt{ПЕРЕСУНЬ: H, (1, 2)} \\
            }
        }
    }
     \\ \\

    \theutterance \stepcounter{utterance}  
    & & & \multicolumn{2}{p{0.3\linewidth}}{
        \cellcolor[rgb]{0.9,0.9,0.9}{
            \makecell[{{p{\linewidth}}}]{
                \texttt{\tiny{[GM$|$GM]}}
                \texttt{ПЕРЕСУНЬ: H, (1, 2)} \\
            }
        }
    }
    & & \\ \\

    \theutterance \stepcounter{utterance}  
    & & & \multicolumn{2}{p{0.3\linewidth}}{
        \cellcolor[rgb]{0.9,0.9,0.9}{
            \makecell[{{p{\linewidth}}}]{
                \texttt{\tiny{[GM$|$GM]}}
                \texttt{Штраф: (1,2) не порожня, а містить '─'.} \\
\\ 
\texttt{Разом ви накопичили 10 з 12 штрафів.} \\
\texttt{Будь ласка, спробуйте ще раз!} \\
            }
        }
    }
    & & \\ \\

    \theutterance \stepcounter{utterance}  
    & & \multicolumn{4}{p{0.6\linewidth}}{
        \cellcolor[rgb]{0.9,0.9,0.9}{
            \makecell[{{p{\linewidth}}}]{
                \texttt{\tiny{[GM$\rangle$P2]}}
                \texttt{Штраф: (1,2) не порожня, а містить '─'.} \\
\\ 
\texttt{Разом ви накопичили 10 з 12 штрафів.} \\
\texttt{Будь ласка, спробуйте ще раз!} \\
            }
        }
    }
    & \\ \\

    \theutterance \stepcounter{utterance}  
    & & & \multicolumn{4}{p{0.6\linewidth}}{
        \cellcolor[rgb]{0.9,0.9,0.9}{
            \makecell[{{p{\linewidth}}}]{
                \texttt{\tiny{[GM$\langle$P2]}}
                \texttt{ПЕРЕСУНЬ: H, (2, 2)} \\
            }
        }
    }
     \\ \\

    \theutterance \stepcounter{utterance}  
    & & & \multicolumn{2}{p{0.3\linewidth}}{
        \cellcolor[rgb]{0.9,0.9,0.9}{
            \makecell[{{p{\linewidth}}}]{
                \texttt{\tiny{[GM$|$GM]}}
                \texttt{ПЕРЕСУНЬ: H, (2, 2)} \\
            }
        }
    }
    & & \\ \\

    \theutterance \stepcounter{utterance}  
    & & & \multicolumn{2}{p{0.3\linewidth}}{
        \cellcolor[rgb]{0.9,0.9,0.9}{
            \makecell[{{p{\linewidth}}}]{
                \texttt{\tiny{[GM$|$GM]}}
                \texttt{Штраф: (2,2) не порожня, а містить '┴'.} \\
\\ 
\texttt{Разом ви накопичили 11 з 12 штрафів.} \\
\texttt{Будь ласка, спробуйте ще раз!} \\
            }
        }
    }
    & & \\ \\

    \theutterance \stepcounter{utterance}  
    & & \multicolumn{4}{p{0.6\linewidth}}{
        \cellcolor[rgb]{0.9,0.9,0.9}{
            \makecell[{{p{\linewidth}}}]{
                \texttt{\tiny{[GM$\rangle$P2]}}
                \texttt{Штраф: (2,2) не порожня, а містить '┴'.} \\
\\ 
\texttt{Разом ви накопичили 11 з 12 штрафів.} \\
\texttt{Будь ласка, спробуйте ще раз!} \\
            }
        }
    }
    & \\ \\

    \theutterance \stepcounter{utterance}  
    & & & \multicolumn{4}{p{0.6\linewidth}}{
        \cellcolor[rgb]{0.9,0.9,0.9}{
            \makecell[{{p{\linewidth}}}]{
                \texttt{\tiny{[GM$\langle$P2]}}
                \texttt{ПЕРЕСУНЬ: H, (1, 3)} \\
            }
        }
    }
     \\ \\

    \theutterance \stepcounter{utterance}  
    & & & \multicolumn{2}{p{0.3\linewidth}}{
        \cellcolor[rgb]{0.9,0.9,0.9}{
            \makecell[{{p{\linewidth}}}]{
                \texttt{\tiny{[GM$|$GM]}}
                \texttt{ПЕРЕСУНЬ: H, (1, 3)} \\
            }
        }
    }
    & & \\ \\

    \theutterance \stepcounter{utterance}  
    & & & \multicolumn{2}{p{0.3\linewidth}}{
        \cellcolor[rgb]{0.9,0.9,0.9}{
            \makecell[{{p{\linewidth}}}]{
                \texttt{\tiny{[GM$|$GM]}}
                \texttt{H moved to (1, 3)} \\
            }
        }
    }
    & & \\ \\

    \theutterance \stepcounter{utterance}  
    & & \multicolumn{4}{p{0.6\linewidth}}{
        \cellcolor[rgb]{0.9,0.9,0.9}{
            \makecell[{{p{\linewidth}}}]{
                \texttt{\tiny{[P1$\langle$GM]}}
                \texttt{Об'єкт 'W' успішно переміщено в (1,3). Ваша оновлена сітка виглядає так:} \\
\texttt{```} \\
\texttt{ 1234567} \\
\texttt{╔═╤═╤═══╗} \\
\texttt{║H│◌│◌◌◌║ 1} \\
\texttt{╟─┴─┼─┬─╢ 2} \\
\texttt{║W◌◌│◌│T║ 3} \\
\texttt{║◌◌◌├─┘◌║ 4} \\
\texttt{║◌┌─┘◌◌◌║ 5} \\
\texttt{╟─┤◌◌◌◌◌║ 6} \\
\texttt{║I│◌C◌◌◌║ 7} \\
\texttt{╚═╧═════╝} \\
\texttt{```} \\
\texttt{Зараз ви граєте раунд 5 з максимум 20.} \\
\texttt{Разом ви накопичили 11 з 12 штрафів.} \\
\texttt{Інший гравець перемістив об'єкт на своїй сітці.} \\
\\ 
\\ 
\texttt{Яка ваша наступна команда?} \\
            }
        }
    }
    & \\ \\

    \theutterance \stepcounter{utterance}  
    & \multicolumn{4}{p{0.6\linewidth}}{
        \cellcolor[rgb]{0.9,0.9,0.9}{
            \makecell[{{p{\linewidth}}}]{
                \texttt{\tiny{[P1$\rangle$GM]}}
                \texttt{ПЕРЕСУНЬ: T, (1, 4)} \\
            }
        }
    }
    & & \\ \\

    \theutterance \stepcounter{utterance}  
    & & & \multicolumn{2}{p{0.3\linewidth}}{
        \cellcolor[rgb]{0.9,0.9,0.9}{
            \makecell[{{p{\linewidth}}}]{
                \texttt{\tiny{[GM$|$GM]}}
                \texttt{ПЕРЕСУНЬ: T, (1, 4)} \\
            }
        }
    }
    & & \\ \\

    \theutterance \stepcounter{utterance}  
    & & & \multicolumn{2}{p{0.3\linewidth}}{
        \cellcolor[rgb]{0.9,0.9,0.9}{
            \makecell[{{p{\linewidth}}}]{
                \texttt{\tiny{[GM$|$GM]}}
                \texttt{T moved to (1, 4)} \\
            }
        }
    }
    & & \\ \\

    \theutterance \stepcounter{utterance}  
    & & \multicolumn{4}{p{0.6\linewidth}}{
        \cellcolor[rgb]{0.9,0.9,0.9}{
            \makecell[{{p{\linewidth}}}]{
                \texttt{\tiny{[GM$\rangle$P2]}}
                \texttt{Об'єкт 'H' успішно переміщено в (1,3). Ваша оновлена сітка виглядає так:} \\
\texttt{```} \\
\texttt{ 1234567} \\
\texttt{╔═╤═╤═══╗} \\
\texttt{║C│◌│◌◌◌║ 1} \\
\texttt{╟─┴─┼─┬─╢ 2} \\
\texttt{║H◌◌│◌│◌║ 3} \\
\texttt{║◌◌◌├─┘◌║ 4} \\
\texttt{║◌┌─┘◌T◌║ 5} \\
\texttt{╟─┤◌W◌◌◌║ 6} \\
\texttt{║◌│◌I◌◌◌║ 7} \\
\texttt{╚═╧═════╝} \\
\texttt{```} \\
\texttt{Зараз ви граєте раунд 6 з максимум 20.} \\
\texttt{Разом ви накопичили 11 з 12 штрафів.} \\
\texttt{Інший гравець перемістив об'єкт на своїй сітці.} \\
\\ 
\\ 
\texttt{Яка ваша наступна команда?} \\
            }
        }
    }
    & \\ \\

    \theutterance \stepcounter{utterance}  
    & & & \multicolumn{4}{p{0.6\linewidth}}{
        \cellcolor[rgb]{0.9,0.9,0.9}{
            \makecell[{{p{\linewidth}}}]{
                \texttt{\tiny{[GM$\langle$P2]}}
                \texttt{ПЕРЕСУНЬ: W, (1, 4)} \\
            }
        }
    }
     \\ \\

    \theutterance \stepcounter{utterance}  
    & & & \multicolumn{2}{p{0.3\linewidth}}{
        \cellcolor[rgb]{0.9,0.9,0.9}{
            \makecell[{{p{\linewidth}}}]{
                \texttt{\tiny{[GM$|$GM]}}
                \texttt{ПЕРЕСУНЬ: W, (1, 4)} \\
            }
        }
    }
    & & \\ \\

    \theutterance \stepcounter{utterance}  
    & & & \multicolumn{2}{p{0.3\linewidth}}{
        \cellcolor[rgb]{0.9,0.9,0.9}{
            \makecell[{{p{\linewidth}}}]{
                \texttt{\tiny{[GM$|$GM]}}
                \texttt{W moved to (1, 4)} \\
            }
        }
    }
    & & \\ \\

    \theutterance \stepcounter{utterance}  
    & & \multicolumn{4}{p{0.6\linewidth}}{
        \cellcolor[rgb]{0.9,0.9,0.9}{
            \makecell[{{p{\linewidth}}}]{
                \texttt{\tiny{[P1$\langle$GM]}}
                \texttt{Об'єкт 'T' успішно переміщено в (1,4). Ваша оновлена сітка виглядає так:} \\
\texttt{```} \\
\texttt{ 1234567} \\
\texttt{╔═╤═╤═══╗} \\
\texttt{║H│◌│◌◌◌║ 1} \\
\texttt{╟─┴─┼─┬─╢ 2} \\
\texttt{║W◌◌│◌│◌║ 3} \\
\texttt{║T◌◌├─┘◌║ 4} \\
\texttt{║◌┌─┘◌◌◌║ 5} \\
\texttt{╟─┤◌◌◌◌◌║ 6} \\
\texttt{║I│◌C◌◌◌║ 7} \\
\texttt{╚═╧═════╝} \\
\texttt{```} \\
\texttt{Зараз ви граєте раунд 6 з максимум 20.} \\
\texttt{Разом ви накопичили 11 з 12 штрафів.} \\
\texttt{Інший гравець перемістив об'єкт на своїй сітці.} \\
\\ 
\\ 
\texttt{Яка ваша наступна команда?} \\
            }
        }
    }
    & \\ \\

    \theutterance \stepcounter{utterance}  
    & \multicolumn{4}{p{0.6\linewidth}}{
        \cellcolor[rgb]{0.9,0.9,0.9}{
            \makecell[{{p{\linewidth}}}]{
                \texttt{\tiny{[P1$\rangle$GM]}}
                \texttt{ПЕРЕСУНЬ: I, (1, 5)} \\
            }
        }
    }
    & & \\ \\

    \theutterance \stepcounter{utterance}  
    & & & \multicolumn{2}{p{0.3\linewidth}}{
        \cellcolor[rgb]{0.9,0.9,0.9}{
            \makecell[{{p{\linewidth}}}]{
                \texttt{\tiny{[GM$|$GM]}}
                \texttt{ПЕРЕСУНЬ: I, (1, 5)} \\
            }
        }
    }
    & & \\ \\

    \theutterance \stepcounter{utterance}  
    & & & \multicolumn{2}{p{0.3\linewidth}}{
        \cellcolor[rgb]{0.9,0.9,0.9}{
            \makecell[{{p{\linewidth}}}]{
                \texttt{\tiny{[GM$|$GM]}}
                \texttt{I moved to (1, 5)} \\
            }
        }
    }
    & & \\ \\

    \theutterance \stepcounter{utterance}  
    & & \multicolumn{4}{p{0.6\linewidth}}{
        \cellcolor[rgb]{0.9,0.9,0.9}{
            \makecell[{{p{\linewidth}}}]{
                \texttt{\tiny{[GM$\rangle$P2]}}
                \texttt{Об'єкт 'W' успішно переміщено в (1,4). Ваша оновлена сітка виглядає так:} \\
\texttt{```} \\
\texttt{ 1234567} \\
\texttt{╔═╤═╤═══╗} \\
\texttt{║C│◌│◌◌◌║ 1} \\
\texttt{╟─┴─┼─┬─╢ 2} \\
\texttt{║H◌◌│◌│◌║ 3} \\
\texttt{║W◌◌├─┘◌║ 4} \\
\texttt{║◌┌─┘◌T◌║ 5} \\
\texttt{╟─┤◌◌◌◌◌║ 6} \\
\texttt{║◌│◌I◌◌◌║ 7} \\
\texttt{╚═╧═════╝} \\
\texttt{```} \\
\texttt{Зараз ви граєте раунд 7 з максимум 20.} \\
\texttt{Разом ви накопичили 11 з 12 штрафів.} \\
\texttt{Інший гравець перемістив об'єкт на своїй сітці.} \\
\\ 
\\ 
\texttt{Яка ваша наступна команда?} \\
            }
        }
    }
    & \\ \\

    \theutterance \stepcounter{utterance}  
    & & & \multicolumn{4}{p{0.6\linewidth}}{
        \cellcolor[rgb]{0.9,0.9,0.9}{
            \makecell[{{p{\linewidth}}}]{
                \texttt{\tiny{[GM$\langle$P2]}}
                \texttt{ПЕРЕСУНЬ: T, (1, 5)} \\
            }
        }
    }
     \\ \\

    \theutterance \stepcounter{utterance}  
    & & & \multicolumn{2}{p{0.3\linewidth}}{
        \cellcolor[rgb]{0.9,0.9,0.9}{
            \makecell[{{p{\linewidth}}}]{
                \texttt{\tiny{[GM$|$GM]}}
                \texttt{ПЕРЕСУНЬ: T, (1, 5)} \\
            }
        }
    }
    & & \\ \\

    \theutterance \stepcounter{utterance}  
    & & & \multicolumn{2}{p{0.3\linewidth}}{
        \cellcolor[rgb]{0.9,0.9,0.9}{
            \makecell[{{p{\linewidth}}}]{
                \texttt{\tiny{[GM$|$GM]}}
                \texttt{T moved to (1, 5)} \\
            }
        }
    }
    & & \\ \\

    \theutterance \stepcounter{utterance}  
    & & \multicolumn{4}{p{0.6\linewidth}}{
        \cellcolor[rgb]{0.9,0.9,0.9}{
            \makecell[{{p{\linewidth}}}]{
                \texttt{\tiny{[P1$\langle$GM]}}
                \texttt{Об'єкт 'I' успішно переміщено в (1,5). Ваша оновлена сітка виглядає так:} \\
\texttt{```} \\
\texttt{ 1234567} \\
\texttt{╔═╤═╤═══╗} \\
\texttt{║H│◌│◌◌◌║ 1} \\
\texttt{╟─┴─┼─┬─╢ 2} \\
\texttt{║W◌◌│◌│◌║ 3} \\
\texttt{║T◌◌├─┘◌║ 4} \\
\texttt{║I┌─┘◌◌◌║ 5} \\
\texttt{╟─┤◌◌◌◌◌║ 6} \\
\texttt{║◌│◌C◌◌◌║ 7} \\
\texttt{╚═╧═════╝} \\
\texttt{```} \\
\texttt{Зараз ви граєте раунд 7 з максимум 20.} \\
\texttt{Разом ви накопичили 11 з 12 штрафів.} \\
\texttt{Інший гравець перемістив об'єкт на своїй сітці.} \\
\\ 
\\ 
\texttt{Яка ваша наступна команда?} \\
            }
        }
    }
    & \\ \\

    \theutterance \stepcounter{utterance}  
    & \multicolumn{4}{p{0.6\linewidth}}{
        \cellcolor[rgb]{0.9,0.9,0.9}{
            \makecell[{{p{\linewidth}}}]{
                \texttt{\tiny{[P1$\rangle$GM]}}
                \texttt{ПЕРЕСУНЬ: C, (1, 1)} \\
            }
        }
    }
    & & \\ \\

    \theutterance \stepcounter{utterance}  
    & & & \multicolumn{2}{p{0.3\linewidth}}{
        \cellcolor[rgb]{0.9,0.9,0.9}{
            \makecell[{{p{\linewidth}}}]{
                \texttt{\tiny{[GM$|$GM]}}
                \texttt{ПЕРЕСУНЬ: C, (1, 1)} \\
            }
        }
    }
    & & \\ \\

    \theutterance \stepcounter{utterance}  
    & & & \multicolumn{2}{p{0.3\linewidth}}{
        \cellcolor[rgb]{0.9,0.9,0.9}{
            \makecell[{{p{\linewidth}}}]{
                \texttt{\tiny{[GM$|$GM]}}
                \texttt{Штраф: (1,1) не порожня, а містить 'H'.} \\
\\ 
\texttt{Разом ви накопичили 12 з 12 штрафів.} \\
\texttt{Будь ласка, спробуйте ще раз!} \\
            }
        }
    }
    & & \\ \\

    \theutterance \stepcounter{utterance}  
    & & \multicolumn{4}{p{0.6\linewidth}}{
        \cellcolor[rgb]{0.9,0.9,0.9}{
            \makecell[{{p{\linewidth}}}]{
                \texttt{\tiny{[P1$\langle$GM]}}
                \texttt{Штраф: (1,1) не порожня, а містить 'H'.} \\
\\ 
\texttt{Разом ви накопичили 12 з 12 штрафів.} \\
\texttt{Будь ласка, спробуйте ще раз!} \\
            }
        }
    }
    & \\ \\

    \theutterance \stepcounter{utterance}  
    & \multicolumn{4}{p{0.6\linewidth}}{
        \cellcolor[rgb]{0.9,0.9,0.9}{
            \makecell[{{p{\linewidth}}}]{
                \texttt{\tiny{[P1$\rangle$GM]}}
                \texttt{СКАЖИ: завершено?} \\
            }
        }
    }
    & & \\ \\

    \theutterance \stepcounter{utterance}  
    & & & \multicolumn{2}{p{0.3\linewidth}}{
        \cellcolor[rgb]{0.9,0.9,0.9}{
            \makecell[{{p{\linewidth}}}]{
                \texttt{\tiny{[GM$|$GM]}}
                \texttt{СКАЖИ: завершено?} \\
            }
        }
    }
    & & \\ \\

    \theutterance \stepcounter{utterance}  
    & & \multicolumn{4}{p{0.6\linewidth}}{
        \cellcolor[rgb]{0.9,0.9,0.9}{
            \makecell[{{p{\linewidth}}}]{
                \texttt{\tiny{[GM$\rangle$P2]}}
                \texttt{Об'єкт 'T' успішно переміщено в (1,5). Ваша оновлена сітка виглядає так:} \\
\texttt{```} \\
\texttt{ 1234567} \\
\texttt{╔═╤═╤═══╗} \\
\texttt{║C│◌│◌◌◌║ 1} \\
\texttt{╟─┴─┼─┬─╢ 2} \\
\texttt{║H◌◌│◌│◌║ 3} \\
\texttt{║W◌◌├─┘◌║ 4} \\
\texttt{║T┌─┘◌◌◌║ 5} \\
\texttt{╟─┤◌◌◌◌◌║ 6} \\
\texttt{║◌│◌I◌◌◌║ 7} \\
\texttt{╚═╧═════╝} \\
\texttt{```} \\
\texttt{Зараз ви граєте раунд 8 з максимум 20.} \\
\texttt{Разом ви накопичили 12 з 12 штрафів.} \\
\texttt{Інший гравець надіслав це повідомлення:} \\
\\ 
\\ 
\texttt{"завершено?"} \\
\\ 
\\ 
\texttt{Яка ваша наступна команда?} \\
            }
        }
    }
    & \\ \\

    \theutterance \stepcounter{utterance}  
    & & & \multicolumn{4}{p{0.6\linewidth}}{
        \cellcolor[rgb]{0.9,0.9,0.9}{
            \makecell[{{p{\linewidth}}}]{
                \texttt{\tiny{[GM$\langle$P2]}}
                \texttt{СКАЖИ: завершено!} \\
            }
        }
    }
     \\ \\

    \theutterance \stepcounter{utterance}  
    & & & \multicolumn{2}{p{0.3\linewidth}}{
        \cellcolor[rgb]{0.9,0.9,0.9}{
            \makecell[{{p{\linewidth}}}]{
                \texttt{\tiny{[GM$|$GM]}}
                \texttt{СКАЖИ: завершено!} \\
            }
        }
    }
    & & \\ \\

    \theutterance \stepcounter{utterance}  
    & & & \multicolumn{2}{p{0.3\linewidth}}{
        \cellcolor[rgb]{0.9,0.9,0.9}{
            \makecell[{{p{\linewidth}}}]{
                \texttt{\tiny{[GM$|$GM]}}
                \texttt{true} \\
            }
        }
    }
    & & \\ \\

    \theutterance \stepcounter{utterance}  
    & & & \multicolumn{2}{p{0.3\linewidth}}{
        \cellcolor[rgb]{0.9,0.9,0.9}{
            \makecell[{{p{\linewidth}}}]{
                \texttt{\tiny{[GM$|$GM]}}
                \texttt{Initial states:} \\
\texttt{```} \\
\texttt{Player 1:} \\
\texttt{ 1234567} \\
\texttt{╔═╤═╤═══╗} \\
\texttt{║◌│◌│◌◌◌║ 1} \\
\texttt{╟─┴─┼─┬─╢ 2} \\
\texttt{║◌◌◌│◌│T║ 3} \\
\texttt{║◌H◌├─┘◌║ 4} \\
\texttt{║◌┌─┘◌◌W║ 5} \\
\texttt{╟─┤◌◌◌◌◌║ 6} \\
\texttt{║I│◌C◌◌◌║ 7} \\
\texttt{╚═╧═════╝} \\
\\ 
\\ 
\texttt{Player 2:} \\
\texttt{ 1234567} \\
\texttt{╔═╤═╤═══╗} \\
\texttt{║◌│◌│◌◌◌║ 1} \\
\texttt{╟─┴─┼─┬─╢ 2} \\
\texttt{║◌◌◌│◌│◌║ 3} \\
\texttt{║◌◌C├─┘◌║ 4} \\
\texttt{║◌┌─┘◌T◌║ 5} \\
\texttt{╟─┤◌W◌◌H║ 6} \\
\texttt{║◌│◌I◌◌◌║ 7} \\
\texttt{╚═╧═════╝} \\
\\ 
\texttt{```} \\
            }
        }
    }
    & & \\ \\

    \theutterance \stepcounter{utterance}  
    & & & \multicolumn{2}{p{0.3\linewidth}}{
        \cellcolor[rgb]{0.9,0.9,0.9}{
            \makecell[{{p{\linewidth}}}]{
                \texttt{\tiny{[GM$|$GM]}}
                \texttt{End states:} \\
\texttt{```} \\
\texttt{Player 1:} \\
\texttt{ 1234567} \\
\texttt{╔═╤═╤═══╗} \\
\texttt{║H│◌│◌◌◌║ 1} \\
\texttt{╟─┴─┼─┬─╢ 2} \\
\texttt{║W◌◌│◌│◌║ 3} \\
\texttt{║T◌◌├─┘◌║ 4} \\
\texttt{║I┌─┘◌◌◌║ 5} \\
\texttt{╟─┤◌◌◌◌◌║ 6} \\
\texttt{║◌│◌C◌◌◌║ 7} \\
\texttt{╚═╧═════╝} \\
\\ 
\\ 
\texttt{Player 2:} \\
\texttt{ 1234567} \\
\texttt{╔═╤═╤═══╗} \\
\texttt{║C│◌│◌◌◌║ 1} \\
\texttt{╟─┴─┼─┬─╢ 2} \\
\texttt{║H◌◌│◌│◌║ 3} \\
\texttt{║W◌◌├─┘◌║ 4} \\
\texttt{║T┌─┘◌◌◌║ 5} \\
\texttt{╟─┤◌◌◌◌◌║ 6} \\
\texttt{║◌│◌I◌◌◌║ 7} \\
\texttt{╚═╧═════╝} \\
\\ 
\texttt{```} \\
            }
        }
    }
    & & \\ \\

    \theutterance \stepcounter{utterance}  
    & & & \multicolumn{2}{p{0.3\linewidth}}{
        \cellcolor[rgb]{0.9,0.9,0.9}{
            \makecell[{{p{\linewidth}}}]{
                \texttt{\tiny{[GM$|$GM]}}
                \texttt{* success: True} \\
\texttt{* lose: False} \\
\texttt{* aborted: False} \\
\texttt{{-}{-}{-}{-}{-}{-}{-}} \\
\texttt{* Shifts: 7.00} \\
\texttt{* Max Shifts: 8.00} \\
\texttt{* Min Shifts: 4.00} \\
\texttt{* End Distance Sum: 14.31} \\
\texttt{* Init Distance Sum: 16.95} \\
\texttt{* Expected Distance Sum: 20.95} \\
\texttt{* Penalties: 12.00} \\
\texttt{* Max Penalties: 12.00} \\
\texttt{* Rounds: 8.00} \\
\texttt{* Max Rounds: 20.00} \\
\texttt{* Object Count: 5.00} \\
            }
        }
    }
    & & \\ \\

\end{supertabular}
}

\end{document}
